
% ════════════════════════════════════════
% 模型检验与分析（06.模型检验.tex）
% ════════════════════════════════════════
\section{模型检验与分析}

为验证本文构建的模型的有效性和稳健性，进行以下检验。本文模型为线性规划优化模型，其最优性由求解器严格保证\cite{boyd2004}，因此检验重点在于约束满足程度、能量守恒精度与参数灵敏度。

\subsection{误差分析}

对问题一的最优解逐项核验其约束满足情况与能量守恒精度，结果如表\ref{tab:error}所示。

\begin{longtable}{>{\centering\arraybackslash}p{0.30\textwidth}>{\centering\arraybackslash}p{0.26\textwidth}>{\raggedright\arraybackslash}p{0.36\textwidth}}
  \caption{模型误差分析结果}
  \label{tab:error} \\
  \toprule
  检验项 & 检验结果 & 说明 \\
  \midrule
  \endfirsthead
  \caption{模型误差分析结果（续）} \\
  \toprule
  检验项 & 检验结果 & 说明 \\
  \midrule
  \endhead
  \bottomrule
  \endfoot
  功率平衡最大残差 & $1.36\times10^{-12}$ & 能量守恒严格满足 \\
  储电量越界次数 & 0 & 全程处于 $[1200,10800]$ kWh \\
  充放电功率越界次数 & 0 & 始终不超过 5000 kW \\
  储电初值约束 & 满足 & $E_0=6000$ kWh，与附录1给定一致 \\
  日循环约束残差 & 0 & $E_0=E_{144}=6000$ kWh \\
\end{longtable}

如表\ref{tab:error}所示，最优解的功率平衡残差处于 $10^{-12}$ 量级，储能荷电状态与充放电功率全程处于约束区间，日循环约束严格成立，验证了模型求解的正确性与数值精度。

\subsection{灵敏度分析}

为检验模型对关键参数的稳健性，分别调整储能充放电效率、最大充放电功率、储电量下限与负荷水平，观察全天购电费的变化，结果如表\ref{tab:sensitivity}所示。

\begin{longtable}{>{\centering\arraybackslash}p{0.30\textwidth}>{\centering\arraybackslash}p{0.30\textwidth}>{\raggedright\arraybackslash}p{0.32\textwidth}}
  \caption{灵敏度分析结果}
  \label{tab:sensitivity} \\
  \toprule
  参数变化 & 输出变化 & 灵敏度评估 \\
  \midrule
  \endfirsthead
  \caption{灵敏度分析结果（续）} \\
  \toprule
  参数变化 & 输出变化 & 灵敏度评估 \\
  \midrule
  \endhead
  \bottomrule
  \endfoot
  $\eta=0.85$ & 36603.41 元 & 费用上升 4.20\% \\
  $\eta=0.95$ & 33767.03 元 & 费用下降 3.87\% \\
  $P_{\max}=4000$ & 35247.87 元 & 费用上升 0.34\% \\
  $P_{\max}=6000$ & 35083.29 元 & 费用下降 0.12\% \\
  $E_{\min}=1500$ & 35329.91 元 & 费用上升 0.58\% \\
  负荷 $+5\%$ & 38834.17 元 & 费用上升 10.55\% \\
  负荷 $-5\%$ & 31552.18 元 & 费用下降 10.18\% \\
\end{longtable}

如表\ref{tab:sensitivity}所示，购电费对充放电效率最为敏感，效率每提高 5\% 可使费用下降约 4\%，说明储能效率是影响经济性的关键因素；充放电功率与储电量下限的小幅变化对费用的影响均不足 1\%，模型稳健；负荷水平按比例影响费用，符合能量平衡的线性规律。综上，模型在各参数摄动下均表现稳定，结果可靠。
